\chapter{APPENDIX: LATEX TRICKS}\label{appendix:latex_tricks}
During this appendix, some general commands are demonstrated.

Using \textit{chapter}, \textit{section}, \textit{subsection} commands, it is possible to add new titles that will be automatically added to the table of contents.

References in the bibliography can be cited using \textit{cite} command. For example \cite{stratton1976needs}. \textit{Cite} command supports also multiple sources. For example: \cite{hirsijarvi2009tutki, virtanen2011esimerkkilahde}. In addition author's name can be referenced using \textit{citeauthor} command. For example: \citeauthor{hirsijarvi2009tutki}.

Using \textit{gls} command, it is easy to reference abbreviations. For example when the command is used first time: \gls{cfd}, abbreviation is fully opened. When the command is used second time, only short version of the abbreviation will be used: \gls{cfd}. This feature is useful, if the text is later refactored. Using \textit{glspl} it is possible to insert plural form of the abbreviation. For example \glspl{orc}.

Titles, figures, tables and many other things can be labeled using \textit{label} command. Later these labels can be automatically referenced using \textit{autoref} command. For example \autoref{appendix:latex_tricks}. Additionally link is automatically created.

Using \textit{itemize} environment and \textit{item} command, it is possible to define unordered lists. For example:
\begin{itemize}[noitemsep]
	\item Item 1
	\item Item 2
\end{itemize}

Using \textit{enumerate} environment and \textit{item} command, it is possible to define ordered lists. For example:
\begin{enumerate}[noitemsep]
	\item Item 1
	\item Item 2
\end{enumerate}

Lists can also be nested naturally. For example:
\begin{enumerate}[noitemsep]
	\item Item 1
	\begin{enumerate}[noitemsep]
		\item Sub item 1
		\item Sub item 2
	\end{enumerate}
\end{enumerate}

When placing figure, LaTeX tries to find optimal placement that use space as efficiently as possible. Sometimes this doesn't end up looking great, because figures are placed too far a way from the text. Figures can be forced to be placed before point x using \textit{FloatBarrier} command.

It is also possible to create new commands using LaTeX using \textit{newcommand}. See example definition and usage of the green check mark bellow:

\newcommand{\greencheckmark}{
	\color{green}\ding{51}\color{black}
}
This is a green check mark: \greencheckmark.

Using appropriate commands, it is possible to create quite advanced structures. For example when presenting survey questions following definitions could be useful:

\makeatletter
\newcommand{\radiobutton}{
	\begin{tikzpicture}
	    % Set radius.
		\pgfmathsetlengthmacro\radius{height("X")/2}
		% Draw circle.
		\draw[radius=\radius] circle;
	\end{tikzpicture}
}
\makeatother

% Create Likert scale description numbers.
\newcommand{\likertscaledescription}{
	1\hspace*{2em}2\hspace*{2em}3\hspace*{2em}4\hspace*{2em}5\hspace*{2em}6
}

% Create Likert scale radio buttons.
\newcommand{\likertscalebuttons}{
	\hspace*{0.2em}\radiobutton\hspace*{1.318em}\radiobutton\hspace*{1.318em}\radiobutton\hspace*{1.318em}\radiobutton\hspace*{1.318em}\radiobutton\hspace*{1.318em}\radiobutton\hspace*{0.2em}
}

% Combine Likert scale components into table.
% Note usage of the parameter usage.
\newcommand{\likertscalesmall}[2]{
	\begin{center}
		\begin{tabular}{c c c}
			& & \\
			& \likertscaledescription & \\ 
			#1 & \likertscalebuttons & #2 \\
			& & \\
		\end{tabular}
	\end{center}
}

\vspace{2em}
\begin{enumerate}
    \item How good this template was?
    \likertscalesmall{Very bad}{Very good}
\end{enumerate}